\documentclass[11pt,a4paper]{article}

\usepackage[utf8]{inputenc}
\usepackage[T1]{fontenc}
\usepackage{lmodern}
\usepackage{amsmath,amssymb,amsfonts}
\usepackage{graphicx}
\usepackage{booktabs}
\usepackage{tabularx}
\usepackage{multirow}
\usepackage{hyperref}
\usepackage{url}
\usepackage{xcolor}
\usepackage{listings}
\usepackage{natbib}
\usepackage{microtype}
\usepackage[margin=1in]{geometry}
\usepackage{enumitem}

\hypersetup{
  colorlinks=true,
  linkcolor=blue,
  citecolor=blue,
  urlcolor=blue,
}

\lstset{
  basicstyle=\ttfamily\small,
  breaklines=true,
  frame=single,
  backgroundcolor=\color{gray!8},
}

% --- Replace before submission ---
\newcommand{\PaperTitle}{Markdown vs.\ Structured Context for LLM Tabular Grounding: A Pilot Study on Ontology Pipeline Artifacts}
\newcommand{\AuthorName}{Your Name}
\newcommand{\AuthorAffil}{ComplianceSpark / Nexcraft}
\newcommand{\AuthorEmail}{you@example.com}
\newcommand{\PaperDate}{May 2026}

\title{\textbf{\PaperTitle}}
\author{
  \textbf{\AuthorName}\\
  \textit{\AuthorAffil}\\
  \texttt{\AuthorEmail}
}
\date{\PaperDate}

\begin{document}
\maketitle

\begin{abstract}
Enterprise ontology pipelines increasingly ground large language models (LLMs) on tabular metadata---column profiles, cardinality classes, and row samples---before enrichment stages such as data-protection classification and relationship inference.
A recurring debate in the semantic-web community asks whether such grounding should use \emph{formal} representations (RDF, JSON-LD, OWL) or \emph{naturalistic} prose (Markdown)~\cite{tonyseale2025mcontology}.
We hold the underlying semantics fixed in a canonical \texttt{TabularContextBundle} and compare only the \textbf{prompt serialization layer}: deterministic Markdown versus compact JSON, across four commercial models (OpenAI, Claude, Gemini, DeepSeek) and four invocation paths (native SDK, LangChain, procedural skill prompt, ontology-foundry provider).
On a pilot benchmark of ten automatically graded factual questions over the \texttt{users\_core} table from a CSOD learning preview dataset, we report accuracy, latency, and qualitative failure modes.
This paper documents the experimental design, open-source harness (\texttt{weekend\_bench.py}), and a reproducible arXiv submission path; \textbf{numeric results should be inserted after running the harness} (placeholders appear as \textbf{TBD} below).
We argue for a \emph{dual serialization} strategy: structured interchange formats for machines and merge semantics; Markdown (or similarly LLM-oriented prose) for in-context grounding when the task is comprehension rather than logical entailment.
\end{abstract}

\noindent\textbf{Keywords:} ontology, knowledge graphs, large language models, tabular profiling, prompt engineering, RDF, Markdown, data governance

\section{Introduction}
\label{sec:intro}

The term \emph{ontology} has become overloaded.
In philosophy, ontology studies what exists; in computer science, Gruber~\cite{gruber1993ontology} defines an ontology as a \emph{formal, explicit specification of a shared conceptualization}.
W3C standards---RDF, RDFS, OWL, SKOS---encode meaning so that reasoners can infer implicit facts (e.g., transitivity of part-of)~\cite{w3crdf11,w3cowl2}.
Practitioners warn of ``McOntology''---labels and shallow taxonomies marketed as ontologies without axioms, shareability, or maintenance discipline~\cite{tonyseale2025mcontology}.

Large language models complicate this picture.
Retrieval-augmented and tool-augmented pipelines~\cite{lewis2020retrieval} inject domain context into prompts rather than into a reasoner.
A popular question on professional networks is whether \emph{informal convergence} through shared foundation models substitutes for \emph{formal agreement} through shared ontologies---or whether formal semantics are simply digested invisibly into the GenAI stack.

This work does not resolve the philosophical debate.
Instead, we isolate an engineering question that arises in production systems such as \textbf{ontology-foundry} and \textbf{ontology-pipeline}:
\begin{quote}
  \emph{Given a fixed tabular semantic artifact, does the choice of in-prompt serialization (Markdown vs.\ JSON/RDF-shaped structure) affect LLM task accuracy, and does that effect depend on model vendor or client framework?}
\end{quote}

Our contributions are:
\begin{enumerate}[leftmargin=*]
  \item A \textbf{semantics-neutral} intermediate representation (\texttt{TabularContextBundle}) and two deterministic serializers (Markdown and JSON) used in enrichment prompts today.
  \item A \textbf{reproducible weekend pilot harness} that runs factual QA with gold answers derived from stored column aggregates, supporting four models and four invocation runtimes.
  \item A \textbf{positioned discussion} on dual serialization for ontology-backed LLM systems, with a roadmap toward RDF/Turtle arms and PII-classification tasks.
\end{enumerate}

\section{Background}
\label{sec:background}

\subsection{Formal ontologies and LLM context}
RDF graphs and OWL axioms excel at interchange, versioning, and machine inference across organizations.
Their cognitive and token costs for \emph{in-context} consumption by autoregressive LLMs are less well measured: triples are sparse, URI-heavy, and may not align with pretraining distributions~\cite{brown2020language}.
Markdown-like tabular summaries, by contrast, mirror technical documentation and data dictionaries found abundantly in pretraining corpora.

These representations are not mutually exclusive.
A robust architecture may export OWL for governance and federation while rendering Markdown (or JSON) \emph{views} for LLM stages---analogous to materialized views over a canonical logical model.

\subsection{TabularContextBundle in ontology-foundry}
The \texttt{ontology-foundry} package profiles database columns (null rate, distinct count, numeric summaries, top-$k$ frequencies) and packages results in a \texttt{TabularContextBundle}:
table identity, optional population metadata, per-column \texttt{ColumnContext} records, and a bounded row sample.
The function \texttt{render\_tabular\_context} produces deterministic Markdown for prompts; no model calls occur during rendering.

\subsection{ontology-pipeline enrichment}
The \texttt{ontology-pipeline} threads the bundle through \texttt{EnrichmentContext} and injects grounding into stages such as:
\begin{itemize}[leftmargin=*]
  \item \textbf{Data protection} --- PII and sensitivity classification with \texttt{aggregates\_only} mode (no raw sample values pre-classification).
  \item \textbf{Column semantics} --- business meaning and semantic units.
  \item \textbf{Relationship and causal inference} --- cross-column and cross-asset hypotheses.
\end{itemize}
Preview runs on a CSOD learning dataset materialize MDL JSON, column aggregates, and samples under \texttt{output/preview/}.

\section{Related work}
\label{sec:related}

\textbf{Structured LLM outputs.} Vendor APIs increasingly support JSON schema constraints~\cite{openai2024structured,anthropic2025claude,google2025gemini}, reducing parse failures in production pipelines.

\textbf{Agent skills and frameworks.} Procedural system prompts (``skills'') specify step-by-step reasoning policies; orchestration frameworks such as LangChain~\cite{langchain2024} wrap heterogeneous providers behind a common \texttt{invoke} interface.
We treat \emph{skill prompts} and \emph{framework wrappers} as experimental factors orthogonal to context format.

\textbf{Grounding evaluation.} Our pilot uses factual QA with gold labels from profiler output---a strict subset of grounding metrics (numeric alignment, span overlap) implemented in \texttt{ontology\_foundry.eval.grounding}.

\section{Method}
\label{sec:method}

\subsection{Design principle: one semantics, many serializations}
All experimental arms share a single \texttt{TabularContextBundle} $B$ built from preview artifacts:
\texttt{users\_core.aggregates.json} and \texttt{users\_core.samples.json}.
Two serializers produce prompt context $C_{\mathrm{md}}$ and $C_{\mathrm{json}}$:
\begin{itemize}[leftmargin=*]
  \item \textbf{Markdown (M):} \texttt{render\_tabular\_context}$(B)$ --- section headers, cardinality interpretation text, embedded JSON sample block.
  \item \textbf{JSON (J):} compact dictionary with parallel numeric fields per column and truncated \texttt{sample\_rows} (no interpretive prose).
\end{itemize}
A scripted validator (planned) will assert numeric parity between $C_{\mathrm{md}}$ and $C_{\mathrm{json}}$ for every reported statistic.

\subsection{Task: factual tabular QA}
Each trial presents $(C, q)$ to model $m$ under runtime $r$, requesting a JSON answer:
\begin{lstlisting}
{"answer": "<value>", "evidence_column": "<column_name>"}
\end{lstlisting}
If the fact is absent from $C$, the model should answer \texttt{UNKNOWN}.
Ten questions (Table~\ref{tab:questions}) target columns in \texttt{users\_core}; gold answers are read directly from aggregate JSON (no hand labeling).

\begin{table}[t]
  \centering
  \caption{Pilot benchmark questions (gold from column aggregates).}
  \label{tab:questions}
  \small
  \begin{tabular}{@{}clcc@{}}
    \toprule
    ID & Question (abbrev.) & Column & Gold \\
    \midrule
    q01 & Distinct count & \texttt{user\_id} & 500 \\
    q02 & Null rate (\%) & \texttt{user\_id} & 0.0 \\
    q03 & Null rate (\%) & \texttt{user\_email} & 23.4 \\
    q04 & Distinct count & \texttt{user\_gender} & 5 \\
    q05 & Cardinality tier & \texttt{user\_gender} & low \\
    q06 & Cardinality tier & \texttt{user\_guid} & identifier \\
    q07 & Distinct count & \texttt{user\_address\_id} & 489 \\
    q08 & Mean & \texttt{user\_appr\_id} & 10.664 \\
    q09 & Maximum & \texttt{user\_appr\_id} & 505 \\
    q10 & Profiled row count & \texttt{applicant\_archived\_flag} & 500 \\
    \bottomrule
  \end{tabular}
\end{table}

\subsection{Models}
\label{sec:models}
Four vendors represent open-weight API, frontier closed, and Google stack:
\begin{center}
  \small
  \begin{tabular}{@{}lll@{}}
    \toprule
    Label & Default model & API key env. \\
    \midrule
    OpenAI    & \texttt{gpt-4o-mini} (override via env.) & \texttt{OPENAI\_API\_KEY} \\
    Claude    & \texttt{claude-sonnet-4-20250514} & \texttt{ANTHROPIC\_API\_KEY} \\
    Gemini    & \texttt{gemini-2.0-flash} & \texttt{GOOGLE\_API\_KEY} \\
    DeepSeek  & \texttt{deepseek-chat} & \texttt{DEEPSEEK\_API\_KEY} \\
    \bottomrule
  \end{tabular}
\end{center}
Temperature is fixed at $0$; structured JSON is requested where the API supports it.

\subsection{Invocation runtimes}
\label{sec:runtimes}
We compare \emph{how} the model is called, holding $(C, q)$ constant:
\begin{description}[leftmargin=2.2cm,style=nextline]
  \item[\textbf{native}] Vendor SDK directly (\texttt{openai}, \texttt{anthropic}, \texttt{google.genai}).
  \item[\textbf{langchain}] \texttt{ChatOpenAI}, \texttt{ChatAnthropic}, \texttt{ChatGoogleGenerativeAI}.
  \item[\textbf{skill}] Native SDK plus a procedural system prompt (\texttt{weekend\_skill\_system.txt}) specifying column lookup steps, percent conversion, and JSON-only replies.
  \item[\textbf{foundry}] \texttt{ontology\_foundry.llm.OpenAIChatProvider} with Pydantic-validated JSON (OpenAI and DeepSeek only).
\end{description}

\subsection{Metrics}
\textbf{Primary:} accuracy --- fraction of questions where parsed \texttt{answer} matches gold under type-specific rules (integer equality, float tolerance, percent with fraction normalization).

\textbf{Secondary:} median latency (ms); optional token counts when APIs expose usage.

\textbf{Planned ablations (post-pilot):} token-matched truncation; Turtle/JSON-LD arms; PII classification F1; multi-table expansion.

\subsection{Experimental matrix}
Let $|\mathcal{M}|=4$ models, $|\mathcal{R}|=4$ runtimes (foundry only on two models), $|\mathcal{F}|=2$ formats, $|\mathcal{Q}|=10$ questions.
The full factorial is up to $4 \times 4 \times 2 \times 10 = 320$ API calls; the harness supports \texttt{--quick} (subset) and \texttt{--resume} for incremental CSV logging.

\section{Implementation}
\label{sec:impl}

\subsection{Software artifacts}
\begin{itemize}[leftmargin=*]
  \item \texttt{ontology-foundry/context/table\_bundle.py} --- bundle schema and Markdown renderer.
  \item \texttt{ontology-pipeline/enrich/grounding.py} --- \texttt{format\_tabular\_grounding} wrapper for enricher prompts.
  \item \texttt{scripts/weekend\_bench.py} --- CLI experiment runner.
  \item \texttt{scripts/weekend\_plot.py} --- accuracy charts for dissemination.
\end{itemize}

\subsection{Results logging}
Each trial appends one CSV row:
\texttt{question\_id, model, runtime, format, correct, answer, gold, latency\_ms, model\_id, error}.
Summaries aggregate by format, model, runtime, and triple interactions.

\section{Results}
\label{sec:results}

\noindent\textbf{Instructions for authors.}
Run the harness after setting API keys, then replace \textbf{TBD} cells in Table~\ref{tab:main} from \texttt{weekend\_bench.py --summarize-only} or the generated plots.
Example command sequence appears in Appendix~\ref{app:repro}.

\begin{table}[t]
  \centering
  \caption{Pilot accuracy (\%) by model and context format, averaged over questions and runtimes unless noted. \textbf{Fill from CSV after experiment.}}
  \label{tab:main}
  \begin{tabular}{@{}lcc@{}}
    \toprule
    Model & Markdown (M) & JSON (J) \\
    \midrule
    OpenAI   & TBD & TBD \\
    Claude   & TBD & TBD \\
    Gemini   & TBD & TBD \\
    DeepSeek & TBD & TBD \\
    \bottomrule
  \end{tabular}
\end{table}

\begin{table}[t]
  \centering
  \caption{Accuracy (\%) by invocation runtime (aggregated over models and formats).}
  \label{tab:runtime}
  \begin{tabular}{@{}lc@{}}
    \toprule
    Runtime & Accuracy \\
    \midrule
    native    & TBD \\
    langchain & TBD \\
    skill     & TBD \\
    foundry   & TBD \\
    \bottomrule
  \end{tabular}
\end{table}

% Uncomment after running weekend_plot.py:
% \begin{figure}[t]
%   \centering
%   \includegraphics[width=0.95\linewidth]{../../results/weekend_TIMESTAMP_overview.png}
%   \caption{Overview dashboard (model, format, runtime breakdown).}
%   \label{fig:overview}
% \end{figure}

\subsection{Hypotheses and interpretation guide}
\textbf{H1 (format):} Markdown exceeds JSON on factual QA when prose interpretation (cardinality tier text) is required (q05--q06).

\textbf{H2 (model):} Frontier models show smaller format gaps than economical APIs.

\textbf{H3 (runtime):} Skill prompts improve JSON arm more than Markdown by compensating for sparse structure.

\textbf{H4 (null result):} If differences fall within run-to-run noise, serialization may be second-order relative to model choice---still a publishable pilot outcome.

Report failure modes anecdotally: conflating fraction vs.\ percent null rates (q03), ignoring \texttt{identifier} tier vocabulary, hallucinating columns absent from $C$.

\section{Discussion}
\label{sec:discussion}

\subsection{Dual serialization}
RDF/OWL remain the right interchange layer for federated ontologies, axiom maintenance, and reasoner-backed QA.
Markdown (or LLM-oriented JSON \emph{views}) are prompt-layer materializations of the same bundle---not replacements for formal semantics.
Production ontology-pipeline deployments already implement this split: MDL and aggregates persist structurally; enrichers consume Markdown grounding.

\subsection{Relation to ``McOntology'' discourse}
Demonstrating measurable differences (or principled null results) under controlled semantics addresses calls for \emph{practical utility showcases} rather than definitional arguments alone.
LinkedIn and technical preprints can link to the open harness as a reproducible artifact.

\subsection{Client frameworks and skills}
LangChain and native SDKs should be benchmarked separately: wrapping overhead and default message templates may shift JSON compliance.
Procedural skill prompts test whether lightweight ``agent policy'' substitutes for format affordances---relevant to Claude-style product features and internal Cursor skills.

\section{Limitations}
\label{sec:limits}
\begin{itemize}[leftmargin=*]
  \item Single table (\texttt{users\_core}), English prompts, one domain (HR/learning).
  \item No RDF/Turtle arm in the pilot (JSON only as structured proxy).
  \item No token-matched comparison; longer Markdown may confound accuracy with budget.
  \item API versions drift; models named here may be superseded.
  \item Automated grading does not capture enrichment quality (PII F1, relationship precision).
\end{itemize}

\section{Future work}
\label{sec:future}
\begin{enumerate}[leftmargin=*]
  \item Add Turtle and JSON-LD serializers; optional OWL reasoner pre-pass arm.
  \item Token-parity truncation and cost-accuracy Pareto curves.
  \item Tasks T1--T3 aligned with \texttt{DataProtectionEnricher} and relationship inference with human gold.
  \item Multi-table CSOD preview and public benchmarks (e.g., TPC-H subsets).
  \item Regression gates via \texttt{ontology\_foundry.eval} for CI release blocking.
\end{enumerate}

\section{Conclusion}
We presented a controlled pilot for comparing Markdown and JSON serializations of a shared \texttt{TabularContextBundle} across four LLM vendors and four invocation runtimes, grounded in ontology-foundry and ontology-pipeline production code.
The benchmark is intentionally small but reproducible; authors insert empirical results after a weekend run.
Whether Markdown wins, JSON wins, or format is neutral, the recommendation is the same: \textbf{keep formal semantics in the bundle and interchange layer; choose prompt serialization empirically for each task and model.}

\section*{Acknowledgments}
Replace with funding, reviewers, and dataset providers.

\section*{Code availability}
\begin{lstlisting}[language=bash]
# Repository paths (monorepo nexcraft):
packages/ontology-foundry/
packages/ontology-pipeline/scripts/weekend_bench.py
packages/ontology-pipeline/scripts/weekend_plot.py
packages/ontology-pipeline/paper/arxiv/
\end{lstlisting}

\bibliographystyle{plainnat}
\bibliography{references}

\appendix
\section{Reproduction checklist}
\label{app:repro}

\begin{lstlisting}[language=bash]
cd packages/ontology-pipeline
python3 -m venv .venv-weekend && source .venv-weekend/bin/activate
pip install -e "../ontology-foundry[llm,langchain,tabular]" \
  -r scripts/requirements-weekend.txt

export OPENAI_API_KEY=...
export ANTHROPIC_API_KEY=...
export GOOGLE_API_KEY=...
export DEEPSEEK_API_KEY=...

python scripts/weekend_bench.py --dry-run
python scripts/weekend_bench.py --quick          # smoke test
python scripts/weekend_bench.py --resume --plot  # full pilot

python scripts/weekend_plot.py --latest --chart all
# Update Table 1-2 in main.tex from printed summary

cd paper/arxiv && make   # builds main.pdf
\end{lstlisting}

\section{Grading rules}
\label{app:grade}
\begin{itemize}[leftmargin=*]
  \item \textbf{integer:} rounded equality.
  \item \textbf{float:} relative tolerance $2\%$ with absolute floor $0.01$.
  \item \textbf{percent:} accept answer as percent (23.4) or fraction (0.234) when question requests percentage.
  \item \textbf{text:} case-insensitive match on cardinality tier labels.
\end{itemize}

\section{Prompt templates}
\label{app:prompts}
\textbf{User template (all runtimes):}
\begin{lstlisting}
Answer ONLY from the context below. If the fact is missing,
set answer to UNKNOWN.
Reply with a single JSON object:
{"answer": "<value>", "evidence_column": "<column>"}

CONTEXT:
<serialized bundle>

QUESTION: <q>
\end{lstlisting}

\textbf{Skill system prompt} (skill runtime only): see \texttt{scripts/weekend\_skill\_system.txt}.

\end{document}
